\chapter{Sentinel Node Biopsy}

\section*{Overview}
Sentinel node biopsy is a procedure that identifies and removes the first lymph node or nodes that drain a cancer, most often in breast cancer and melanoma, to stage the disease.

\section*{Alternative Names}
Sentinel lymph node biopsy; SLNB; sentinel node mapping

\section*{Principle}
A tracer (radioactive colloid, blue dye, or both) is injected near the tumor and travels to the sentinel node, which represents the status of the whole nodal basin. Removing it predicts whether cancer has spread.

\section*{History}
The sentinel node concept was introduced for melanoma by Morton in 1992 and rapidly extended to breast cancer, replacing routine axillary dissection for staging.

\section*{Why It Is Done}
Performed for early breast cancer and melanoma without clinically involved nodes, to stage the axilla or groin accurately with minimal surgery.

\section*{Is This a Primary or Secondary Test or Procedure}
Primary staging procedure for node-negative-appearing solid tumors.

\section*{How It Is Performed}
A radioactive tracer or dye is injected before or during surgery, and a gamma probe and visual inspection identify the sentinel node, which is removed and sent for pathology.

\section*{Preparation}
Injection of tracer, and discussion of the implications of a positive result.

\section*{Contraindications and Precautions}
Clinically positive nodes and prior surgery in the draining basin may preclude the procedure.

\section*{Risks}
Seroma, infection, allergic reaction to dye, and a small risk of lymphedema.

\section*{Aftercare and Recovery}
Pathology determines whether further node dissection or axillary radiation is needed.

\section*{Outcome and Follow-up}
A negative sentinel node predicts a very low chance of cancer in remaining nodes; a positive node changes treatment.

\section*{Expected Results}
No cancer in the sentinel lymph node.

\section*{Follow-up}
Further treatment or surveillance based on nodal status.

\section*{When to Seek Medical Care}
Follow the procedure team's specific instructions. Seek urgent medical assessment for severe or worsening pain, uncontrolled bleeding, fever or signs of infection, trouble breathing, chest pain, fainting, new weakness or confusion, inability to urinate, or any rapidly worsening symptom. The appropriate warning signs depend on the procedure and the patient's medical condition.

\section*{References}
\begin{enumerate}
  \item MedlinePlus. Sentinel node biopsy. U.S. National Library of Medicine; 2025.
  \item Current Medical Diagnosis and Treatment. McGraw-Hill; 2025.
\end{enumerate}

\clearpage
